\section{Sistemi Lineari e Rouché-Capelli}

\begin{concept}{Richiami di Teoria: Esistenza delle Soluzioni}
Un sistema in forma matriciale è $Ax = b$.
\textbf{Teorema di Rouché-Capelli:} Il sistema è \textit{compatibile} (ammette soluzioni) se e solo se il rango della matrice incompleta è uguale a quello della matrice completa:
$$ \text{rk}(A) = \text{rk}(A|b) = r $$
Se questa condizione è verificata:
\begin{itemize}
    \item Se $r = n$ (incognite), la soluzione è \textbf{unica} ($\infty^0$).
    \item Se $r < n$ (incognite), esistono \textbf{infinite soluzioni} dipendenti da $n - r$ variabili libere (o gradi di libertà).
\end{itemize}
\end{concept}

\begin{logicbox}
\textbf{Pattern 2: Discutere un Sistema Parametrico}\\
Se è presente un parametro $k$ o $a$:
1. Scrivi la matrice $(A|b)$ associata al sistema.
2. Procedi subito alla riduzione verso la forma a \textit{scalini} (non applicare formule complesse a priori).
3. Studia i pivot: 
   \begin{itemize}
       \item Trova per quali $k$ i pivot (gli elementi in diagonale o nei gradini) si annullano (es. se pivot $= k-2$, pongo $k=2$).
       \item Se il pivot si annulla e hai l'ineguaglianza "zero = costante non nulla", il sistema è \textbf{Impossibile}.
       \item Per i valori in cui nessun pivot cruciale si annulla, trovi una e \textbf{una sola soluzione}.
       \item Per i valori in cui un'intera riga si annulla a zero (sia matrice che termini noti), hai \textbf{infinite} soluzioni e un rango più basso.
   \end{itemize}
\end{logicbox}

\begin{exercisebox}[title={Esercizio 1: Risoluzione Standard}]
Risolvere il sistema lineare omogeneo associato alla matrice:
$$ A = \begin{pmatrix} 1 & -2 & 3 \\ 2 & -4 & 6 \\ -1 & 2 & -3 \end{pmatrix} $$
\end{exercisebox}

\begin{solbox}
Costruiamo la matrice per il sistema omogeneo ($b=0$):
$$ \left(\begin{array}{ccc|c} 1 & -2 & 3 & 0 \\ 2 & -4 & 6 & 0 \\ -1 & 2 & -3 & 0 \end{array}\right) $$
Procediamo con l'eliminazione di Gauss: \\
$R_2 \to R_2 - 2R_1$: $2 - 2(1) = 0, -4 - 2(-2) = 0, 6 - 2(3) = 0$. Tutta una riga di zeri.\\
$R_3 \to R_3 + R_1$: $-1 + 1 = 0, 2 + (-2) = 0, -3 + 3 = 0$. Tutta una riga di zeri.\\
La matrice ridotta diventa:
$$ \left(\begin{array}{ccc|c} 1 & -2 & 3 & 0 \\ 0 & 0 & 0 & 0 \\ 0 & 0 & 0 & 0 \end{array}\right) $$
Il rango è $r=1$ e le incognite sono $n=3$. Abbiamo $\infty^{3-1} = \infty^2$ soluzioni (2 parametri liberi).
Le colonne senza pivot sono la $y$ e la $z$. Poniamo $y=t, z=s$.
L'unica equazione rimasta è $x - 2y + 3z = 0 \implies x = 2t - 3s$.
Le combinazioni di soluzioni assumono la forma:
$$ S = \{ (2t-3s, t, s) \vert t,s \in \mathbb{R} \} $$
\end{solbox}
